\exercisesection{Regular rings}

\begin{enumerate}
\def\labelenumi{\arabic{enumi})}
\item
  Let A be a Noetherian local ring. For A to be regular, it is necessary and sufficient that we have \(\mathrm{di}_{\mathrm{A}}(\kappa_{\mathrm{A}})<+\infty\) .
\item
  Let A be a regular ring and p a prime ideal of A. Prove that the completion of A with respect to the p-adic topology is an integral domain.
\item
  Let A be a ring, and \(\mathbf{T} = (T_i)_{i\in I}\) a finite family of indeterminates. Prove that the following conditions are equivalent:
\end{enumerate}

\begin{enumerate}
\def\labelenumi{(\roman{enumi})}
\item
  the ring A is regular;
\item
  the ring A{[}T{]} is regular;
\item
  the ring A{[}{[}T{]}{]} is regular.
\end{enumerate}

\begin{enumerate}
\def\labelenumi{\arabic{enumi})}
\setcounter{enumi}{3}
\item
  Let A be a Noetherian ring of dimension \(\leqslant 2\) Show that for A to be regular, it is necessary and sufficient that the dual of every finitely generated A-module be projective (i.e., \(\mathfrak{p}\) a prime ideal of \(\Lambda\) ; if the second condition is satisfied, first prove that \(\mathrm{A}_{\mathfrak{p}}\) is regular when \(\operatorname{prof}(A_{\mathfrak{p}}) \leqslant 1\) and then that \(\mathrm{dp}_{\mathrm{A}_{\mathfrak{p}}}(\kappa(\mathfrak{p})) = 2\) when \(\operatorname{prof}(A_{\mathfrak{p}}) = 2\) ).
\item
  Let A be a Noetherian local ring of finite homological dimension; deduce from exerc. 16 of A, X, p.~208 another proof of the fact that A is regular (observe that one has dim \(_{\kappa_A}\) \(\mathfrak{m}_A/\mathfrak{m}_A^2 \leqslant \mathrm{dh}(A) \leqslant \mathrm{prof}(A)\) ).
\item
  Let A be a regular local ring, \((x_{1},\ldots,x_{n})\) and \((y_{1},\ldots,y_{n})\) maximal intersecting sequences of elements of \(m_{A}\) , generating ideals x and y respectively. We assume that we have \(y \subset x\) , so that there exist elements \(a_{ij}\) of A such that we have \(y_{i} = \sum_{j} a_{ij} x_{j}\) for \(1 \leqslant i \leqslant n\) . Prove that the class of \(\det(a_{ij})\) in A/y does not depend on the matrix \((a_{ij})\) chosen, and that its annihilator is the image of x; in particular, one has \(\det(a_{ij}) \neq 0\) (prove that the canonical homomorphism \(\operatorname{Ext}_{\mathrm{A}}^{n}(\mathrm{A}/x,\mathrm{A}) \to \operatorname{Ext}_{\mathrm{A}}^{n}(\Lambda/\mathfrak{y},\Lambda)\) is injective, then show using Koszul complexes that it identifies with the homomorphism A/x → A/y induced by the homothety of ratio \(\det(a_{ij})\) ).
\end{enumerate}

Let A be a regular local ring of dimension 3, and let f be a nonzero element of \(m_{A}\) Show that the following conditions are equivalent:

\begin{enumerate}
\def\labelenumi{(\roman{enumi})}
\item
  the ring A/(f) is not factorial;
\item
  there exists an integer \(n \geqslant 2\) and a matrix \(X \in \mathbf{M}_n(\mathrm{A})\) with elements in \(\mathfrak{m}_{\mathrm{A}}\) , such that \(f = \det(\mathrm{X})\) .
\end{enumerate}

(Under hypothesis (ii), prove that the ideal of A/(f) generated by the classes of the minors \(X^{11}, \ldots, X^{1n}\) of X, where \(X^{1p}\) is obtained by deleting the first row and the column of index p, has height 1 but is not principal. Under hypothesis (i), let p be a prime ideal of height 2 of A whose image in \(\Lambda/(f)\) is not principal; deduce from exerc. 11 of §3 that p is generated by the minors of order n-1 of a matrix \(Y \in \mathbf{M}_{n,n-1}(\mathbf{A})\) with elements in \(m_{\Lambda}\) . Write f as the determinant of a matrix \(\widetilde{Y}\) obtained by adding a column \(Y_{n}\) to Y; if an element of \(Y_{n}\) is invertible, reduce to the case where \(Y_{n}\) belongs to the canonical basis of \(\Lambda^{n}\) .

¶ 8) Let A be a regular local ring, \(\mathbf{x} = (x_1, \ldots, x_d)\) a coordinate system of A, \(\rho : A \to B\) an injective ring homomorphism making B a finitely generated A-module. Prove that the following conditions are equivalent:

\begin{enumerate}
\def\labelenumi{(\roman{enumi})}
\item
  for every integer \(p\) , the element \(\rho(x_1 \ldots x_d)^p\) does not belong to the ideal of B generated by \(\rho(x_1^{p+1}), \ldots, \rho(x_d^{p+1})\) ;
\item
  the A-module \(\rho(A)\) is a direct factor of B.
\end{enumerate}

These conditions are satisfied if A is a Q-algebra (exerc. 3 of § 2).

To prove (ii)⇒(i), adapt the proof of loc. cit. For the converse implication, reduce to the case where A is complete, and set, for \(p \geqslant 0\) \(A_{p} = A/(x_{1}^{p+1}, \ldots, x_{d}^{p+1})\) and \(B_{p} = B/(x_{1}^{p+1}, \ldots, x_{d}^{p+1})\) B. Observe that the class dc \((x_{1} \ldots x_{d})^{p}\) generates the socle of \(A_{p}\) and deduce from this that the homomorphism \(\rho_{p}: A_{p} \to B_{p}\) induced by \(\rho\) is injective. Prove that \(\rho_{p}\) admits a retraction \(A_{p}\) -linear. Conclude using E, III, p.~60, example II.

\begin{enumerate}
\def\labelenumi{\arabic{enumi})}
\setcounter{enumi}{8}
\tightlist
\item
  Let \(\rho: A \to B\) a local homomorphism of Noetherian local rings, \(N\) a finitely generated Macaulay B-module. Suppose that the ring \(A\) is regular and that we have \(\dim_{B}(N) = \dim(A) + \dim_{B/\mathfrak{m}_{A}B}(N/\mathfrak{m}_{A}N)\) . Prove that the A-module \(N\) is flat (argue by induction on \(\dim(A)\) ; if \(x \in \mathfrak{m}_{A} - \mathfrak{m}_{A}^{2}\) we will show that \(x\) is \(N\) -regular and that the homomorphism A/xA → B/xB and the (B/xB)-module N/xN satisfy the same hypotheses as ρ and N).
\end{enumerate}

¶ 10) Let A be a Noetherian ring, M a finitely generated A-module. For every integer \(n \geqslant 0\) , we denote by \(X_{n}\) the set of elements \(\mathfrak{p}\) of \(\operatorname{Spec}(A)\) such that \(\dim_{A_{\mathfrak{p}}}(M_{\mathfrak{p}}) - \operatorname{prof}_{A_{\mathfrak{p}}}(M_{\mathfrak{p}}) \geqslant n\) . Let \(k\) an integer \(\geqslant 0\) .

\begin{enumerate}
\def\labelenumi{\alph{enumi})}
\item
  For M to satisfy property \(S_{k}\) (§ 1, exerc. 8), it is necessary and sufficient that one have codim( \(X_{n}\) , Supp(M)) \(\geqslant n + k\) for every \(n \geqslant 0\) .
\item
  Assume that \(X_{n}\) is closed in \(\operatorname{Spec}(A)\) for every \(n \geqslant 0\) . Let p be an element of \(\operatorname{Supp}(M)\) . For the \(A_{p}\) -module \(M_{p}\) satisfies the property \(S_{k}\) (§ 1, exerc. 8), it is necessary and sufficient that one have \(\operatorname{codim}(X_{n}^{\prime}, \operatorname{Supp}(M)) \geqslant n + k\) for every \(n \geqslant 0\) and every component \(X_{n}^{\prime}\) of \(X_{n}\) containing p.
\item
  Assume that the ring A is presentable. Deduce from b) that the set of prime ideals p of A such that \(M_{p}\) satisfies the property \(S_{k}\) is open in Spec(A).
\end{enumerate}

\begin{enumerate}
\def\labelenumi{\arabic{enumi})}
\setcounter{enumi}{10}
\tightlist
\item
  Let A be the ring \(\mathbf{Z}[\mathbf{X}]\) , where \(\mathbf{X} = (X_i)_{i\in I}\) is a finite family of indeterminates. Let J be an ideal of A generated by monomials.
\end{enumerate}

\begin{enumerate}
\def\labelenumi{\alph{enumi})}
\item
  Let \(i \in I\) ; let \(\mathrm{P}_{1}, \ldots, \mathrm{P}_{r}\) ; \(\mathrm{Q}_{1}, \ldots, \mathrm{Q}_{s}\) of monomials generating J, where the \(\mathrm{P}_{j}\) are divisible by \(\mathrm{X}_{i}\) while c the \(\mathrm{Q}_{k}\) are not. Prove that the transporter J: \(\mathrm{X}_{i}\) (ideal of A consisting of the elements P such that \(\mathrm{X}_{i} \mathrm{P} \in \mathrm{J}\) ) is generated by the monomials \(\mathrm{X}_{i}^{-1} \mathrm{P}_{1}, \ldots, \mathrm{X}_{i}^{-1} \mathrm{P}_{r}\) ; \(\mathrm{Q}_{1}, \ldots, \mathrm{Q}_{s}\) .
\item
  Therefore deduce that if \(J: X_i\) is equal to \(J\) or \(A\) for every \(i \in I\) , \(J\) is the ideal generated by the variables \(X_i\) that belong to \(J\) .
\item
  Let M be an A-module, and let p be an integer. Assume that \(\operatorname{Tor}_{p}^{A}(A/J',M)\) (resp. \(\operatorname{Ext}_{A}^{p}(A/J',M)\) ) is zero for every ideal \(J'\) containing J and generated by some of the \(X_{i}\) . Prove that \(\operatorname{Tor}_{p}^{A}(A/J,M)\) (resp. \(\operatorname{Ext}_{A}^{p}(A/J,M)\) ) is zero (let \(J'\) an ideal containing J, generated by monomials, such that \(\operatorname{Tor}_{p}^{A}(A/J',M) \neq 0\) , and maximal for these properties; deduce from the exact sequence
\[
0 \rightarrow \mathrm{A} / (\mathrm{J} ^ {\prime}: \mathrm{X} _ {i}) \longrightarrow \mathrm{A} / \mathrm{J} ^ {\prime} \longrightarrow \mathrm{A} / (\mathrm{J} ^ {\prime} + \mathrm{AX} _ {i}) \rightarrow 0
\]

and hence b) that J' is generated by the elements \(X_{i}\) that it contains, which contradicts the hypothesis).

\item
  Hence deduce the inequality \(\mathrm{dp}_{\mathrm{A}}(\mathrm{A}/\mathrm{J}) \leqslant \mathrm{Card}(\mathrm{I})\) .

\item
  Let A be a ring, and \(\mathbf{x} = (x_{i})_{i \in I}\) a finite family of elements of A; suppose that for every subset I' of I the family \((x_{i})_{i \in I'}\) is completely secant for A. Let J be an ideal of A generated by elements \(x^{\alpha}\) , where \(\alpha\) ranges over a finite subset F of \(N^{1}\) . Prove the inequality dp \(_{A}\) (A/J) ≤ Card(I) (let J \(_{0}\) be the ideal of the ring A \(_{0}\) = Z{[}(X \(_{i}\) ) \(_{i \in I}\) {]} generated by the \(X^{\alpha}\) for \(\alpha \in F\) ; deduce from c) that Tor \(_{p}^{A_{0}}\) (A \(_{0}\) /J \(_{0}\) , A) is zero, and conclude by considering a projective resolution of the A \(_{0}\) -module A \(_{0}\) /J \(_{0}\) of length ≤ n).
\end{enumerate}

¶ 12) Let A be a regular local ring, M a finitely generated A-module, such that the A-module End \(_{A}\) (M) let it be free. We propose to prove that the following conditions are equivalent:

\begin{enumerate}
\def\labelenumi{(\roman{enumi})}
\item
  the A-module M is free;
\item
  the A-module M is reflexive;
\item
  \(\operatorname{Ext}_{\mathrm{A}}^{1}(\mathbf{M},\mathbf{M}) = 0\)
\end{enumerate}

\begin{enumerate}
\def\labelenumi{\alph{enumi})}
\item
  Prove the implication (iii) \(\Rightarrow\) (ii) (use exerc. 2 of VII, § 4).
\item
  Prove the converse implication (use Exerc. 4 if dim(A) ≤ 2, then reason by induction on dim(A) using Exerc. 18 of § 1).
\item
  Assume \(\dim(A) \leqslant 3\) Show the implication (ii) \(\Rightarrow\) (i) (use exerc. 4 if \(\dim(A) \leqslant 2\) ; if \(\dim(A) = 3\) , observe that (ii) implies \(\mathrm{dp}_{A}(M) \leqslant 1\) , then that (iii) and remark 2 of \(n^{\circ} 3\) imply that M is projective).
\item
  Assume M is reflexive. Let x be an element of \(m_{A} = m_{A}^{2}\) ; denote by B the ring A/xA, by N the B-module M/xM and by \(N^{*}\) its dual. Using b), prove that the B-module \(\operatorname{End}_{\mathbb{B}}(\mathbb{N})\) is free, then that \(\operatorname{End}_{\mathbb{B}}(\mathbb{N}^{*})\) is free (deduce from VII, § 4, no. 2 that \(\operatorname{End}_{\mathbb{B}}(\mathbb{N}^{*})\) is identified with the bidual of \(\operatorname{End}_{\mathbb{B}}(\mathbb{N})\) ).
\item
  Prove that (ii) implies (i) (argue by induction on the dimension of A, assumed ≥ 4; with the notation of d), the induction hypothesis implies that the B-module N* is free; deduce from this, using exerc. 18 of § 1 and prop. 7 of no. 4, that one has prof \(_{A}\) (Ext \(^{1}_{A}\) (M, B)) ≥ 1, then that Ext \(^{1}_{A}\) (M, A), which is of finite length by the induction hypothesis, is zero; conclude that the B-module M*/xM* is identified with N*, hence is free, then that the A-module M* is free and finally that M is free.)
\item
  Deduce from these results another proof of the fact that the ring A is factorial.
\end{enumerate}

\begin{enumerate}
\def\labelenumi{\arabic{enumi})}
\setcounter{enumi}{11}
\tightlist
\item
  Let A be a ring, and \((x_{1},\ldots,x_{n})\) a family of elements of A, generating an ideal I. We say that the family \((x_{1},\ldots,x_{n})\) is A-independent if the classes of \(x_{i}\) in I/I \(^{2}\) form a basis of this A/l-module.
\end{enumerate}

\begin{enumerate}
\def\labelenumi{\alph{enumi})}
\item
  Every completely intersecting family for A is A-independent.
\item
  Let \((x_{1},\ldots,x_{n})\) a A-independent family, with \(x_{1}=x_{1}^{\prime}x_{1}^{\prime\prime}\) . Prove that the families \((x_{1}^{\prime},\ldots,x_{n})\) and \((x_{1}^{\prime\prime},\ldots,x_{n})\) are A-independent, and that one has
\end{enumerate}

\[
\lg_ {\mathrm{A}} \left(\mathrm{A} / \left(x _ {1}, \dots , x _ {n}\right)\right) = \lg_ {\mathrm{A}} \left(\mathrm{A} / \left(x _ {1} ^ {\prime}, \dots , x _ {n}\right)\right) + \lg_ {\mathrm{A}} \left(\mathrm{A} / \left(x _ {1} ^ {\prime \prime}, \dots , x _ {n}\right)\right)
\]

(observe that the quotient \((x_{1}^{\prime},\ldots,x_{n})/(x_{1},\ldots,x_{n})\) is isomorphic to \(A/(x_{1}^{\prime\prime},\ldots,x_{n})\) ).

¶ 13) Let \(p\) be a prime number, \(q\) a power of \(p\) Let A be a reduced Noetherian local ring of characteristic \(p\) . One denotes by \(\mathrm{A}^q\) the subring of A consisting of the elements \(a^q\) for \(a \in \mathrm{A}\) . Prove that for A to be regular, it is necessary and sufficient that the \(\mathrm{A}^q\) -module A be flat.

(Reduce to the case where A is complete, hence a quotient of a formal power series algebra \(B = \kappa_{A}[[X_{1}, \ldots, X_{n}]]\) by an ideal b, which may be assumed to be contained in \(m_{B}^{2}\) . If A is regular, then b = 0. If A is flat over \(A^{q}\) let \(x_{i}\) be the image of \(X_{i}\) in A; prove that the family \((x_{1}^{q}, \ldots, x_{n}^{q})\) is A-independent (exercise 12), and deduce from this the equality \(\lg_{\mathrm{A}}(\mathrm{A}/(x_{1}^{q}, \ldots, x_{n}^{q})) = q^{n}\) . Show that this implies \(b \subset (X_{1}^{q}, \ldots, X_{n}^{q})\) ; prove similarly that one has \(b \subset (X_{1}^{q^{s}}, \ldots, X_{n}^{q^{s}})\) for every s, whence finally b = 0.)

